\section{Euler screening before window geometry and shared disc observers}
\label{sec:window-euler}

The kernel update removes wide matching elimination, but repeated source
geometry still costs time. This continuation implements the Euler-corner
criterion preserved from report 81 in
Section~\ref{sec:planar-overlay-refinements}. It screens updated windows
before building their standard matching geometry. Positive-Euler windows
retain their ordinary standard-ray and independently replayed disc queries.
A second change shares validated producer geometry with those disc queries,
while leaving independent certificate reconstruction in place.

\subsection{Compile Euler characteristic from the actual source}

For each global edge choose one local representative. A normal surface's
intersection weight there is
\[
 w_{\{a,b\}}(x)=x_a+x_b+\sum_{q=0}^2x_{4+q}
                                  -x_{4+Q(a,b)},
\]
where $Q(a,b)$ is the quadrilateral type disjoint from that edge.
Matching makes this weight independent of the representative. Count normal
arcs on every boundary face and one representative of every paired face;
on face $f$ their count is $\sum_{v\ne f}x_v+\sum_qx_{4+q}$.
Each elementary normal disc supplies one face of the induced surface cell
structure. Thus
\[
 \chi(x)=\sum_{\text{global edges}}w_e(x)
       -\sum_{\text{global faces}}\operatorname{arcs}_f(x)
       +\sum_{t,j}x_{t,j}=\beta\cdot x.
\]
The native compiler forms the integer vector $\beta$ by precisely these
three sums. It neither expands discs nor assumes that a local edge is a
separate global edge. Independent native coordinate analysis remains the
reference for the resulting Euler values.

Let $P$ be the retained all-Q source forest potentials of
Section~\ref{sec:window-basis}. The signed standard extension is
$U(q)=(Pq,q)$. Compile once the coefficients
\[
 \gamma_j=\beta_{Q,j}+\sum_c\beta_{\triangle,c}P_{c,j},
 \qquad \lambda_v=\sum_{c\in G_v}\beta_{\triangle,c}.
\]
The second sum evaluates $\beta$ on the actual vertex link. The compact
manifold source has $\lambda_v\in\{1,2\}$; the compiler checks these values.
For an updated basis $K'$ and parameters $z$, the existing projection gives
corner forms $\ell_c(z)=(PK')_c z$. Canonical normalization subtracts the
minimum in each global vertex group. Consequently
\[
 \chi(\operatorname{lift}(q))=\gamma\cdot q
                   -\sum_v\lambda_v\min_{c\in G_v}\ell_c(z),
 \qquad q=\sum_a z_aK'_a.
\]
This evaluates the exact Euler value without allocating the full coordinate
vector or reconstructing matching and edge-weight geometry for that vector.
It includes every physical corner, even those that the later projected
kernel would contract or omit as a pure link factor.

\subsection{A complete nonpositive screen and a conservative positive branch}

The displayed formula is positively homogeneous and convex: the minima of
linear forms are concave and all link coefficients are positive. On the
compact normalized nonnegative Q section, a convex combination of corners
has Euler value at most the same combination of their values. Nonpositive
values at every corner therefore exclude positive Euler everywhere in the
sector. An essential normal disc is canonical and has nonzero Q coordinates
and Euler characteristic one, so it cannot occur in an excluded sector.
Positive Euler alone still proves neither connectedness nor essentiality.

The implementation uses the existing exact normalized chart and polygon
clipping at matching nullity at most three. Empty, point and segment sections
remain possible even when the raw nullity is larger. It evaluates every
corner needed for a nonpositive answer; a callback interruption propagates
and cannot supply an exclusion. As soon as a positive corner is found,
ordinary standard discovery resumes. The projected forms are reused when
building the standard kernel, so the screen does not project the source a
second time. At nullity above three it continues through the existing
complete generic arrangement/support paths without an Euler exclusion.

Nullity one has a cheaper branch. The canonical basis row has a positive
unit free coordinate; any negative entry therefore makes its nonnegative
cone empty. This check precedes corner projection. Otherwise the sole
primitive direction's Euler sign can be evaluated directly from scalar
corner forms. No normalized rational rescaling or affine chart is needed.
The sign, rather than an unnormalized numerical maximum, is the criterion.

These exclusions are internal search decisions, not portable negative knot
certificates. The original radius-one/two window, its candidate order and
its positive proof format are retained. A completed nonpositive sector can
be skipped, while any positive-Euler sector still receives the full ordinary
standard-ray query. Finite shared work and orbit caps remain inconclusive.
The radius-zero recognizer path stays unchanged.

\subsection{Reuse producer geometry without weakening source replay}

The public disc-count interface continues to validate the actual source,
coordinates and all options. It then delegates to a private observer on that
validated geometry. The window already owns a read-only prepared source and
has freshly analysed each candidate vector against it; it can call the same
private observer directly. This removes an additional source preparation
and coordinate analysis from each positive-Euler candidate query.

The canonical disc core is also compared exactly with the analysed input
rows. If they agree, its geometry is already available. Otherwise the reduced
core is analysed afresh. Equality handles genuine primitive canonical rays
without assuming that every supplied normal vector has those properties.
Orbit-census fallbacks retain their established behavior.

No certificate verifier uses this private observer. Independent disc replay
still reconstructs the source and input geometry, removes links and scale
arithmetically, and checks the core proof separately. Successful diagram
proofs still authenticate the canonical diagram exterior. The reused producer
records are not an accepted substitute for any of these checks.

\subsection{Cost and scope}

The source Euler compilation costs $O(T+\operatorname{nnz}P)$ integer
additions after the planner has supplied $P$. For a $d$-row updated basis,
projection costs $O(d\operatorname{nnz}P_S)$ arithmetic operations.
The implementation scans physical corner forms at each tested Q corner;
with $m$ such corners its Euler work is
$O(m d(T+k))$ plus the exact normalized-chart construction.
It does not claim the fully aggregated $O(T+k^2)$ precheck proposed in
report 81. Eliminating the remaining repeated physical-corner scans is a
separate possible improvement.

Excluded sectors avoid standard matrix allocation, minimum-envelope/overlay
construction, full ray lifts and their topology queries. Positive sectors
pay for screening before their unchanged complete standard search.
Producer geometry reuse removes repeated validation rather than changing
compressed observer mathematics. Bit costs, coefficient growth, output and
independent replay remain charged. The preceding local quasi-polynomial
bound for logarithmic nullity/radius survives; no universal centre family,
distance bound or general quasi-polynomial unknot algorithm follows.

\subsection{Validation and measured behavior}
All 1,407 maintained tests pass in 240.354 seconds; the fifteen focused
tests pass in 18.920 seconds. Every tested small source window compares
the screen with complete freshly reconstructed standard rays. Both the
compiled linear functional and canonical formula agree with independent
native coordinate analysis, including scaling by $2^{90}+1$.
A genuine nullity-four source continues through the existing complete
fallback. The private prepared observer reproduces all seven genuine
fixture queries and their independently replayed disc proofs.

The 85-source audit exactly preserves all disabled legacy outputs, all
enabled verdicts and all 31 native positives. Seven ordinary window discs
receive fresh Regina confirmation. Work-capped cases fall from thirty to
twenty-eight, while complete window misses rise from twenty-four to
twenty-six. The newly completed sources are \code{random-31} and
\code{random-144}, at 1,552,225 and 1,751,788 guard calls. They still yield
inconclusive native knot results. The genus-one positive keeps its 32
sector queries and 31 basis updates, emits eighteen rather than twenty-two
rays, and prunes twenty sectors, including empty Q sections. Its work drops
from 127,635 to 109,610 calls. All 598 captured source pins remain fixed.
The enabled incumbent corpus observations are frozen preceding-release data;
disabled comparisons and all updated queries are performed afresh.
The 234.535-second audit is not a paired timing against earlier audits.

An initial Euler-only pilot completes 120 measured calls and forty warmups.
It shows only small gains on two window cases and approximately parity on
the trefoil. Its raw evidence and the exact source snapshots needed to
recover all 598 original pins are preserved. It does not establish a
uniform speedup. Profiling then motivates the prepared-observer integration.

The final comparison freezes five preceding-release modules at
\code{1539f353a}: disc counting, window basis construction, residual queries,
seed queries and recognition. Their local imports are redirected only to
corresponding frozen dependencies; no timed patch context is added.
Five serial interleaved four-arm rounds complete 200 measured calls and
forty warmups. Native timings include the whole configured window query,
independent successful source replay and serialized output. Both arms
finish all configured work. Positive certificate digests agree exactly.

\begin{samepage}
\noindent\textbf{Complete native queries with radius-two windows.}
\begin{center}
\small
\begin{tabular}{lrrrrr}
\toprule
Input & previous ms & current ms & old/new & old A/A & new A/A\\
\midrule
empty circle & 6.030 & 5.927 & 0.985 & 0.983 & 1.004\\
optimized positive & 59.189 & 58.038 & 1.015 & 1.009 & 0.982\\
genus one miss & 281.501 & 246.270 & 1.143 & 0.989 & 1.067\\
trefoil & 545.240 & 422.983 & 1.233 & 1.044 & 1.014\\
figure eight & 994.493 & 754.190 & 1.312 & 1.007 & 0.981\\
\bottomrule
\end{tabular}
\end{center}

\end{samepage}

The first genus-one seed comparison has a new-arm A/A ratio of 1.067.
A separate fifteen-round repeat completes sixty measured calls and four
warmups with the same full query, replay and serialization scope. Its
paired old/new gain is 1.228, with A/A controls 1.020 and 1.030.
Both runs are retained; the noisier first row is not the sole evidence
for a genus-one native speedup.

\begin{samepage}
\noindent\textbf{Complete recognition with radius-two windows enabled.}
\begin{center}
\small
\begin{tabular}{lrrrrr}
\toprule
Input & previous ms & current ms & old/new & old A/A & new A/A\\
\midrule
empty circle & 0.019 & 0.018 & 1.020 & 0.973 & 0.958\\
optimized positive & 50.866 & 50.965 & 1.002 & 0.999 & 0.993\\
genus one miss & 299.243 & 239.963 & 1.230 & 0.993 & 0.983\\
trefoil & 561.791 & 447.279 & 1.241 & 0.999 & 1.013\\
figure eight & 1040.929 & 791.445 & 1.318 & 0.990 & 0.999\\
\bottomrule
\end{tabular}
\end{center}

\end{samepage}

Complete recognition improves about 1.230 times on the genus-one case,
1.241 on the trefoil and 1.318 on the figure-eight. The latter two include
the established complete fallback after native misses. These rows have
A/A controls within about two percent of parity. Early positive controls
remain approximately unchanged. The empty-recognition row is a microsecond
control; its ratio is not interpreted as a kernel speedup. These gains are
relative to the preceding enabled-window implementation. Radius zero remains
much faster on the small controls, so the window stage stays off by default.

\begin{sloppypar}
Runtime release \code{f95a86fdf} contains the implementation, regression
checks and \code{normal\_orbit\_research.window\_euler} driver. Use
\code{audit --fresh-regina} or \code{benchmark --rounds 5} from \code{fast/}.
Publication checks compare 598 frozen/current runtime pins, 595 preceding
release evidence pins, all timing outcomes and both saved exploratory/repeat
records. They independently replay all 31 saved positive proofs with Euler
compilation, prepared observers, matching updates, residual planning,
coherent extraction, ray enumeration and orbit producers disabled.
The checked \code{openai/math} head remains \code{fd4aeeb2ee4f}; no new
external revision supplies an additional complexity claim.
\end{sloppypar}

